%!TeX root=../thesis.tex

\chapter{Background}
\label{cap:background}
\label{chap:background}
\label{ch:background}

This chapter establishes the theoretical foundations for the empirical investigation
carried out in this work. It begins with an overview of software architecture and
architectural uncertainty, then introduces hypotheses engineering and the ArchHypo framework in technical details, describes Internal
Developer Portals and their role as platforms for enabling tools such as Hypo-Stage,
presents the Hypo-Stage tool itself, and concludes with a summary of findings from the systematic mapping study that mapped the gaps between architectural derivation methods and practical tool support \citep{souza2019}.

% ---------------------------------------------------------------------------
\section{Software Architecture and Architectural Uncertainty}
\label{sec:sw-arch}

Software architecture refers to the high-level structure of a software system, encompassing its components, their relationships, and the design principles that guide
its evolution ~\citep{bass2021}. However, in organizations that maintain software at the core of their business developed and operated by multiple teams,
the ever-happening architectural decisions shape not only technical quality attributes such as performance, reliability, and maintainability of systems, 
but also the ability of engineering teams to work independently and deliver changes quickly and efficiently at a speed that needs to match the speed of the business ~\citep{skelton2019}.

More embracing to that rapidly changing business scenarios, \citet{richardsFord2020} added that the design decisions themselves are part of the software architecture itself, 
making it not only a static structure or frozen blueprint of a system, but a living body of reasoning that is constantly being refined and adapted to the changing business needs. 
This framing is adopted in this work to ground the concept of Architectural Uncertainty: the rationale behind any architectural choice is valid only relative to the information available at the time of the decision, 
and that information is always incomplete and subject to change.

In other words, Architectural Uncertainty arises whenever teams must make structural decisions under
incomplete information. This is a normal condition in continuous-delivery environments
where business requirements, user behavior, and technology options are only partially
known. When uncertainty is left implicit, teams either defer decisions indefinitely
(accumulating technical debt) or make premature actions that may constrain the
architecture evolution of their systems. Both outcomes can slow down delivery and degrade system
quality over time, turning the software architecture into what Foote and Yoder \citeyearpar{footeYoder1999} famously characterized as a ``Big Ball of Mud.''

Moreover, the concept of ``fast flow'', popularized through Team Topologies \citep{skelton2019}, emphasizes that organizations should design both their
teams and their architecture to enable a continuous, rapid stream of software changes.
Achieving fast flow presupposes that teams can make local architectural decisions
without waiting for a central architect to approve them. This in turn demands that architectural
uncertainties be made explicit and that lightweight mechanisms exist for teams to assess, prioritize, and revisit them as evidence accumulates. 
This is where Hypotheses Engineering may be used to play a key role in the architecture work, as proposed by \citet{silva2024tse}.

% ---------------------------------------------------------------------------
\section{Hypotheses Engineering and ArchHypo}
\label{sec:archhypo}

Hypotheses engineering is an approach to managing technical uncertainty by formulating
architectural decisions as testable hypotheses rather than settled facts. Instead of
asserting that a particular technology or pattern \emph{will} satisfy a quality
requirement, teams state that it \emph{might} do so and then design experiments to
gather evidence to validate or invalidate the hypothesis.

The ArchHypo technique, proposed by \citet{silva2024tse}, operationalized this idea for
software architecture. Grounded in Design Science Research (DSR)
and evolved from a body of patterns for agile architecture first disseminated at
conferences and through corporate training reported in~\citep{silva2024ieeesw}, ArchHypo provides
a structured vocabulary and a set of practices that fit within iterative development
workflows where architecture uncertainty is a constant challenge for rapidly-evolving software systems.

% ------- 2.2.1
\subsection{Core Concepts}
\label{subsec:archhypo-core}

An \emph{architectural hypothesis} is a falsifiable statement that captures an uncertainty relevant to the design of a software architecture. 
It is important to distinguish the concept from adjacent concepts:

\begin{itemize}
  \item \textbf{Hypothesis versus\ Risk:} A risk involves a known event with estimable
  probability and impact. An architectural hypothesis, by contrast, stems from a \emph{lack of information} sufficient to make a decision. 
  The distinguishing feature is the epistemic gap, not the probability of an outcome as observed in \citet{silva2024tse}.

  \item \textbf{Hypothesis versus\ Requirement:} A requirement specifies what the system
  must do or how well it must do it. A hypothesis may be triggered by a requirement
  (e.g., a quality-attribute goal that cannot yet be confirmed as achievable), but the
  hypothesis captures the uncertainty about the \emph{path} to satisfying that
  requirement, not the requirement itself.

  \item \textbf{Hypothesis versus\ Architectural Decision.} An architectural decision
  resolves a design choice. A hypothesis \emph{precedes and informs} decisions: it
  identifies what must be learned or validated before a decision can responsibly be
  made (~\citep{silva2024tse}).
\end{itemize}

\begin{table}[htbp]
\centering
\caption{Distinguishing architectural hypotheses from adjacent concepts \citep{silva2024tse}.}
\label{tab:hypothesis-concepts}
\begin{tabular}{p{0.22\linewidth} p{0.35\linewidth} p{0.35\linewidth}}
\toprule
\textbf{Concept} & \textbf{Definition} & \textbf{Relation to Architectural Hypothesis} \\
\midrule
Risk & A known event with estimable probability and impact. & A hypothesis stems from an \emph{epistemic gap} --- lack of information --- not from a known probability. \\
\midrule
Requirement & Specifies what the system must do or how well it must do it. & A requirement may \emph{trigger} a hypothesis when the path to satisfying it is uncertain; they are not the same artifact. \\
\midrule
Architectural Decision & Resolves a design choice. & A hypothesis \emph{precedes and informs} a decision; it captures what must be learned before a decision can responsibly be made. \\
\bottomrule
\end{tabular}
\end{table}

Hypotheses arise from two principal sources:

\begin{itemize}
  \item (a)~\emph{requirements}: when quality attribute targets are poorly specified or the team lacks information about whether a given approach can satisfy them; and 
  \item (b)~\emph{solutions}: when the consequences of an existing or candidate architectural solution are unknown.
\end{itemize}

% ------- 2.2.2
\subsection{Assessment}
\label{subsec:archhypo-assessment}

Each hypothesis is evaluated along two independent dimensions:

\begin{enumerate}
  \item \textbf{Uncertainty level.} How far the team is from being able to refute or
  accept the hypothesis. Crucially, uncertainty level is \emph{not} equivalent to
  probability. A hypothesis that is almost certainly true and a hypothesis that is almost
  certainly false both represent \emph{low} uncertainty, because in either case the team
  already has enough information to act. High uncertainty means the team genuinely
  does not know which way the evidence will fall.

  \item \textbf{Impact level.} The potential consequences --— in terms of rework,
  schedule, cost, or quality degradation --— if the hypothesis is invalidated. A high
  uncertainty combined with a high impact signals that a hypothesis demands urgent
  attention.
\end{enumerate}

Both dimensions can be assessed using a 3-point or 5-point qualitative scale;
assessments are made independently and relatively, not on absolute probability as reported in \citet{silva2024tse}.

% ------- 2.2.3
\subsection{Technical Plans}
\label{subsec:archhypo-technical-plans}

For hypotheses with significant uncertainty or impact, teams formulate a
\emph{technical plan}: a set of concrete actions whose goal is to (a)~accept or refute
the hypothesis, (b)~reduce its uncertainty level, (c)~reduce its potential impact, or
(d)~define criteria under which the decision can responsibly be
deferred. ArchHypo identifies the following plan strategies as observed in \citet{silva2024ieeesw}:

\begin{itemize}
  \item \textbf{Architectural Spike.} A time-boxed, isolated experiment to explore a
  technology or design option without coupling the result to the production codebase.
  The spike produces information rather than deliverable functionality.

  \item \textbf{Tracer Bullet.} A thin, end-to-end implementation within the actual
  system, exercising all architectural layers to validate an integration or performance
  characteristic under realistic conditions.

  \item \textbf{Software Analytics.} Using monitoring, profiling, or usage data already
  collected from the running system to gather evidence about a hypothesis without
  additional development effort.

  \item \textbf{Modularity\,/\,Flexibility.} Restructuring the affected subsystem to
  reduce coupling or increase replaceability, thereby lowering the \emph{impact} of the
  hypothesis rather than resolving its uncertainty. This strategy is useful when full
  resolution is premature but protection against a bad outcome is warranted.

  \item \textbf{Guideline Implementation.} Tailoring the development process—for
  example, adopting a coding standard, a review protocol, or a testing policy—so that
  the team systematically generates evidence about the hypothesis as a by-product of
  normal work. This was the most frequently used strategy in the evaluation reported by \citet{silva2024tse}.

  \item \textbf{Architectural Trigger.} A condition-based deferral in which the team
  agrees to act on the hypothesis only when a specific observable event occurs (e.g.,
  when the number of concurrent users exceeds a threshold, or when a performance
  regression appears in monitoring). Until the trigger fires, the hypothesis is
  deliberately deferred.

  \item \textbf{Agile Landing Zones.} Adopting architectural patterns that explicitly
  preserve flexibility, creating ``landing zones'' into which future decisions can be
  inserted without major disruption, regardless of which direction the evidence
  eventually points.
\end{itemize}

\begin{table}[htbp]
\centering
\caption{ArchHypo technical plan strategies and their primary effect \citep{silva2024ieeesw,silva2024tse}.}
\label{tab:technical-plans}
\begin{tabular}{p{0.25\linewidth} p{0.35\linewidth} p{0.30\linewidth}}
\toprule
\textbf{Strategy} & \textbf{Description} & \textbf{Primary Effect} \\
\midrule
Architectural Spike & Time-boxed isolated experiment; produces information, not deliverable functionality. & Reduces uncertainty. \\
\midrule
Tracer Bullet & Thin end-to-end implementation in the real system to validate integration or performance. & Reduces uncertainty under realistic conditions. \\
\midrule
Software Analytics & Uses existing monitoring or profiling data to gather evidence. & Reduces uncertainty without extra development effort. \\
\midrule
Modularity / Flexibility & Restructures the subsystem to reduce coupling or increase replaceability. & Reduces impact (not uncertainty). \\
\midrule
Guideline Implementation & Tailors the development process (coding standards, review protocols) to generate evidence as a by-product of normal work. & Reduces uncertainty incrementally; most frequently used \citep{silva2024tse}. \\
\midrule
Architectural Trigger & Defers action until a specific observable condition is met. & Manages timing of decision; reduces premature cost. \\
\midrule
Agile Landing Zones & Adopts patterns that preserve flexibility for future insertions. & Reduces impact proactively. \\
\bottomrule
\end{tabular}
\end{table}

% ------- 2.2.4
\subsection{Lifecycle}
\label{subsec:archhypo-lifecycle}

Hypotheses progress through a defined lifecycle: they begin in an \emph{open} state
and transition to either \emph{validated} or \emph{invalidated} as evidence from
technical plans accumulates. The \emph{responsible moment} is the key decision point in the hypothesis lifecycle:
it is the point at which the team has accumulated enough evidence to act on the
architectural decision --- neither so early that the decision would be premature, nor
so late that deferral would create unnecessary risk.

% ------- 2.2.5
\subsection{Empirical Evaluations}
\label{subsec:archhypo-evaluations}

\paragraph{IEEE Software evaluation (Catch Solve startup).}
The first evaluation of ArchHypo in a real operating environment was reported in an
IEEE Software article in \citet{silva2024ieeesw}. The setting was Catch
Solve, a startup offering web testing and monitoring services based in Italy's South
Tyrol region. Four architectural hypotheses were formulated and
tracked over approximately ten months, covering quality attributes of availability,
reusability of test templates, usability, and scalability of virtual machines. For
example, the hypothesis \emph{``The lack of redundancy can cause problems with
application availability for the customers''} was assessed with a low impact and very
low uncertainty, which led the team to defer investment in redundant infrastructure and
instead set up an Architectural Trigger based on the number of new
customers. The study showed that naming and tracking
uncertainties explicitly helped the technical lead decide when to invest in architectural
improvements versus when to wait, and validated the overall structure of the ArchHypo
approach. A key insight from this evaluation was that explicit hypothesis management
changed how the team prioritized technical work: rather than treating architectural
uncertainty as urgent and tacit, it became visible, assessed, and deliberately managed
within a defined lifecycle.

\paragraph{TSE evaluation (PropSEP at Jetsoft).}
The larger and more controlled evaluation, published in the IEEE Transactions on
Software Engineering ~\citep{silva2024tse}, examined the application of ArchHypo within
Jetsoft, a Brazilian software development company, on the PropSEP project---the budget
proposal system of the Planning and Management Secretariat of the State of São Paulo.
PropSEP is mission-critical: the approved budget it processes amounts to approximately
317~billion~BRL ~\citep{silva2024tse}. The evaluation ran across six sprints of seven
business days each (May--August 2022) and involved eight participants (P1--P8), all
holding computer science degrees with certifications in Oracle, WebSphere, or GeneXus.
Ten architectural hypotheses were identified, spanning six quality attributes
(performance, security, reliability, flexibility, usability, and team productivity), and
seven distinct types of architectural practice were employed in the resulting technical
plans~\citep{silva2024tse}.

The study produced nine key findings:

\begin{enumerate}
  \item Some hypotheses carried uncertainty that was monitored throughout the entire
  project duration without reaching full resolution, demonstrating that the technique
  accommodates long-horizon uncertainties ~\citep{silva2024tse}.

  \item ArchHypo enabled the team to avoid upfront architectural design, distributing
  structural decisions across iterations as evidence became available ~\citep{silva2024tse}.

  \item Participants reported improvements to the development process in terms of
  safety, transparency, and predictability ~\citep{silva2024tse}.

  \item Hypotheses spanning different quality attributes were handled using qualitatively
  different technical strategies, illustrating the creativity the technique
  affords~\citep{silva2024tse}.

  \item Guideline Implementation (process tailoring) emerged as the most frequently
  employed strategy across all technical plans~\citep{silva2024tse}.

  \item The team collectively recognized that managing hypotheses improved the quality
  of their architectural decision-making ~\citep{silva2024tse}.

  \item Post-release evidence pointed to a positive impact on product quality: no issues
  related to the managed hypotheses were detected after deployment ~\citep{silva2024tse}.

  \item \textbf{Changes in the team's way of working constituted a significant adoption
  obstacle.} Mixing architectural and hypothesis-related tasks with functional
  development created estimation difficulties, and participants described the additional
  work as not feeling directly connected to product delivery ~\citep{silva2024tse}.

  \item \textbf{All eight participants unanimously considered the technique hard to
  learn.} Specific difficulties included: mapping architectural concerns to the hypothesis
  format, formulating appropriate technical plan actions, articulating the assumptions
  behind a hypothesis, and understanding hypotheses as the primary driver of
  architectural action. No participant felt prepared to apply ArchHypo without
  continued expert support after the initial workshop ~\citep{silva2024tse}. 
  One participant explicitly stated that *''more training would make the team less dependent on the team of architects''* (P3) ~\citep{silva2024tse}. 
  The authors themselves acknowledge this structural dependency in their discussion. 
  The absence of a supporting tool that encodes this knowledge into the workflow risks ArchHypo to remain a technique that only functions when an experienced architect is actively steering its application.
\end{enumerate}

Findings ~8 and ~9 directly motivate the tooling contribution of this work: a tool that
embeds ArchHypo scaffolding in the daily engineering workflow, reducing the cognitive
overhead of manual application and knowledge transfer.

\begin{table}[htbp]
\centering
\caption{Summary of prior empirical evaluations of ArchHypo.}
\label{tab:archhypo-evaluations}
\begin{tabular}{p{0.18\linewidth} p{0.18\linewidth} p{0.18\linewidth} p{0.18\linewidth} p{0.18\linewidth}}
\toprule
\textbf{Study} & \textbf{Setting} & \textbf{Scope} & \textbf{Key Finding} & \textbf{Limitation} \\
\midrule
\citep{silva2024ieeesw} & Catch Solve startup (Italy) & 4 hypotheses, 10 months & Explicit hypothesis tracking changed how team prioritized architectural investments. & Single practitioner; no tool support. \\
\midrule
\citep{silva2024tse} & PropSEP / Jetsoft (Brazil) & 10 hypotheses, 6 sprints, 8 participants & ArchHypo distributed architectural work across iterations; all participants found technique hard to learn. & Manual application; expert steering required; no supporting tool. \\
\midrule
This work & Large-scale tech organization (Latin America) & Multiple teams, scaled product & To be reported --- first evaluation in a tool-mediated setting via Hypo-Stage. & Single organizational context; Backstage dependency. \\
\bottomrule
\end{tabular}
\end{table}

% ---------------------------------------------------------------------------
\section{Internal Developer Portals and Backstage}
\label{sec:idp-backstage}

The adoption of microservice architectures and the pursuit of team autonomy through
practices such as Team Topologies have created significant complexity in the daily
workflow of software engineering teams. In response, the industry has converged around the
concept of \emph{Internal Developer Portals} (IDPs) as a platform-engineering
discipline aimed at reducing this complexity while preserving team
independence ~\citep{correa2025}.

% ------- 2.3.1
\subsection{The Role of Internal Developer Portals (IDPs)}
\label{subsec:idp-role}

An Internal Developer Portal serves as the primary interface between developers and the
underlying platform infrastructure. Unlike a simple wiki or a project-management tool,
an IDP is an operational hub that aggregates tools, services, documentation, and data
into a unified ``pane of glass'' ~\citep{correa2025}. Its primary goal is to reduce
\emph{cognitive load}: by abstracting the complexity of infrastructure and standardizing
workflows, IDPs enable what practitioners call \emph{Golden Paths} -- recommended,
supported, and automated paths for building and deploying software ~\citep{correa2025}.

This goal aligns directly with the theory of autonomous stream-aligned teams summarized in Team Topologies guidance \citep{skelton2019}: if teams are to own full slices of a product
without depending on a central gatekeeper, the platform must lower the barriers to
acting on that ownership. An IDP operationalizes this by ensuring that the information
and tooling a team needs—service catalog, documentation, CI/CD pipelines, monitoring
dashboards, and domain-specific plugins—are accessible from a single environment,
without context-switching.

Conceptually, IDPs sit at the intersection of two ideas that are highly relevant to this
work. First, they embody the principle of \emph{decentralized architecture}: rather
than requiring teams to escalate architectural concerns to a central architecture group,
a well-designed IDP makes it possible to encode architectural guidance, constraints, and
tools directly into the developer workflow. Second, they are the natural delivery vehicle
for hypothesis-engineering support: because ArchHypo must become part of the daily
engineering cadence --- not a separate off-sprint activity ---embedding it in an IDP portal where
developers already spend their working time is strategically motivating.

% ------- 2.3.2
\subsection{Backstage: Open Platform for Internal Developer Portals}
\label{subsec:backstage}

Backstage\footnote{\url{https://github.com/backstage/backstage}} is an open-source framework for building Internal Developer Portals,
originally developed at Spotify to manage the complexity arising from rapid growth and
an expanding microservice estate. Spotify's internal version of Backstage serves as the
central nervous system for its infrastructure, used by over 2,700 engineers to manage
more than 14,000 software components~\citep{backstage2024}. Spotify open-sourced the
framework in March~2020, and it was subsequently donated to the Cloud Native
Computing Foundation (CNCF) as an Incubating project ~\citep{backstage2024}. As of
2025, more than 3,000 companies have adopted Backstage for their own
portals ~\citep{backstage2024}.

Backstage is not a rigid application but rather a modular ecosystem driven by
configuration. Its architecture rests on three layers: (1)~the \emph{Core Framework},
which manages essential utilities including the plugin loader, routing, authentication,
and error handling; (2)~the \emph{Plugin Ecosystem}, built around a split-plugin model
where frontend plugins handle views and UI components while backend plugins manage
API services, data processing, and storage; and (3)~\emph{External Integrations}, through
which backend plugins serve as gateways to external systems and APIs, allowing the
portal to aggregate data from cloud providers, CI/CD tools, and other
infrastructure ~\citep{correa2025}.

\begin{table}[htbp]
\centering
\caption{Backstage architecture layers and their role \citep{correa2025}.}
\label{tab:backstage-layers}
\begin{tabular}{p{0.25\linewidth} p{0.35\linewidth} p{0.32\linewidth}}
\toprule
\textbf{Layer} & \textbf{Responsibilities} & \textbf{Relevance to Hypo-Stage} \\
\midrule
Core Framework & Plugin loader, routing, authentication, error handling. & Provides the lifecycle hooks that Hypo-Stage's frontend and backend plugins register into. \\
\midrule
Plugin Ecosystem & Frontend plugins (views, UI); backend plugins (API services, data, storage). & Hypo-Stage is implemented as a split plugin: the frontend renders hypotheses and assessments; the backend persists data and exposes the API. \\
\midrule
External Integrations & Backend plugins as gateways to external systems (cloud, CI/CD, APIs). & Hypo-Stage integrates with the Backstage catalog to auto-associate hypotheses with software entities. \\
\bottomrule
\end{tabular}
\end{table}

Crucially, the plugin architecture allows organizations to extend the portal's
functionality by integrating custom tools directly into the developer's daily workflow,
without forcing context-switching. It is precisely this extensibility that makes Backstage
the appropriate delivery vehicle for the Hypo-Stage tool: a plugin can introduce a new class of
first-class object -- the architectural hypothesis -- into the same environment where a
developer already inspects service health, reviews documentation, and triggers CI/CD
runs \citep{correa2025}.

\citet{correa2025} demonstrated the technical feasibility of building such a plugin
and validated the integration approach. That work established the initial version of
the Hypo-Stage plugin for Backstage, which constitutes the artifact under empirical evaluation in this work.

% ---------------------------------------------------------------------------
\section{Hypo-Stage: The Initial Tool}
\label{sec:hypo-stage}

Hypo-Stage is an open-source Backstage plugin developed
as an undergraduate capstone project by Pedro Henrique Mariano Corrêa, supervised by
the author of this work and co-supervised by Prof.\ Paulo
Meirelles \citep{correa2025}. Its purpose is to operationalize ArchHypo within the Spotify's IDP Backstage, embedding hypothesis management into the engineering platform workflow.

The initial version of Hypo-Stage, corresponding to the code up to commit
\texttt{ceee509776508081ebdd473c2c4f710b8ef55947} in the \texttt{hypo-stage}
repository,\footnote{Source repository: \url{https://github.com/ArchHypo/hypo-stage}. The baseline revision for the initial tool design referenced throughout this thesis is commit \texttt{ceee509776508081ebdd473c2c4f710b8ef55947}.} provides the
following capabilities:

\begin{itemize}
  \item \textbf{Hypothesis management}: creating, editing, listing, and deleting
  architectural hypotheses, each linked to a Backstage catalog entity (e.g., a
  service, component, or system). Integration with the Backstage catalog means that
  hypotheses are automatically associated with the software entities they affect, without
  requiring separate registration.

  \item \textbf{Uncertainty and impact assessment}: rating each hypothesis on 5-point
  Likert scales for both uncertainty level and impact level, with free-text rationale
  fields. The use of a 5-point scale rather than a binary or 3-point scale provides
  finer-grained prioritization while remaining tractable for teams without formal risk
  management backgrounds.

  \item \textbf{Technical planning}: attaching one or more technical plan items to a
  hypothesis, specifying action type (spike, tracer bullet, trigger, etc.), description,
  target date, expected outcome, and documentation links.

  \item \textbf{Visualization}: a chart plotting hypotheses on an
  uncertainty-vs-impact plane to support prioritization, making the relative urgency of
  hypotheses visible at a glance.

  \item \textbf{Status lifecycle}: tracking each hypothesis through the defined lifecycle
  states (open~$\rightarrow$~validated or invalidated), preserving the history of status
  transitions to support retrospective analysis.
\end{itemize}

By integrating with Backstage, Hypo-Stage leverages the portal's service catalog so
that hypotheses are automatically associated with the software entities they affect. This
design reduces context-switching: developers manage architectural hypotheses in the
same environment where they consult documentation, inspect service health, and access
CI/CD pipelines.

\paragraph{Known limitations:} The initial tool design assumed familiarity with the
ArchHypo vocabulary, which poses a barrier for teams new to hypothesis engineering.
The case study \citep{correa2025} reconstructed the Catch Solve startup scenario
digitally but did not involve real-time use by a team. Furthermore, the tool lacked
features for collaborative assessment (e.g., multiple raters per hypothesis) and for
tracking the evolution of assessments over time. These gaps also motivated the empirical
evaluation and refinement effort promoted in this work.

\begin{table}[htbp]
\centering
\caption{Hypo-Stage initial capabilities and known limitations identified prior to this study \citep{correa2025}.}
\label{tab:hypostage-capabilities}
\begin{tabular}{p{0.30\linewidth} p{0.32\linewidth} p{0.30\linewidth}}
\toprule
\textbf{Capability} & \textbf{Description} & \textbf{Known Limitation} \\
\midrule
Hypothesis management & Create, edit, list, delete hypotheses linked to Backstage catalog entities. & Assumed prior ArchHypo vocabulary familiarity; no onboarding guidance. \\
\midrule
Uncertainty \& impact assessment & 5-point Likert scales per dimension with free-text rationale. & No collaborative multi-rater assessment; evolution of assessments not tracked over time. \\
\midrule
Technical planning & Attach plan items with action type, description, target date, and outcome. & Action types not enforced or explained in-tool; required expert facilitation to use correctly. \\
\midrule
Visualization & Uncertainty-vs-impact scatter plot for prioritization. & Static; not updated reactively as assessments changed. \\
\midrule
Status lifecycle & Tracks hypotheses through open $\rightarrow$ validated / invalidated states. & No history of intermediate state transitions preserved. \\
\bottomrule
\end{tabular}
\end{table}

% ---------------------------------------------------------------------------
\section{Deriving Architecture from Requirements: The Tooling Gap}
\label{sec:tooling-gap}

The systematic mapping study \citep{souza2019} examined
the existing body of work on methods for deriving architectural models from requirements
specifications. Starting from 2,575 candidate papers, the authors applied a PICOC
(Population, Intervention, Comparison, Outcome, Context) search structure and a quality
assessment checklist -- composed of general criteria (weighted at 25\%) and
method-specific criteria (weighted at 75\%) -- to select 39 high-quality primary studies
for detailed analysis. Their analysis employed the NIMSAD (Normative Information Model-based Systems Analysis and Design) framework to characterize the
structure of each reviewed method.

Among their key findings, these were directly relevant to this work:

\begin{enumerate}
  \item \textbf{Dependence on tacit knowledge.} Every reviewed method relied, to some
  degree, on the architect's experience and intuition. No method provided a fully
  automated or guided derivation path accessible to less-experienced practitioners~\citep{souza2019}.

  \item \textbf{Limited support for less-experienced practitioners.} Methods assumed a
  skilled architect capable of navigating trade-offs among quality attributes. Teams without senior architectural expertise had little actionable guidance~\citep{souza2019}.

  \item \textbf{Weak or absent tool support.} Although some tools were identified—most
  notably QuaDAI and Monarch—they addressed only fragments of the derivation process
  and offered only partial coverage of the architectural reasoning cycle. The mapping
  found no comprehensive tool integrating hypothesis formulation, assessment, and
  tracking within a team's daily workflow~\citep{souza2019}.

  \item \textbf{Insufficient evaluation.} Many proposed methods lacked rigorous
  empirical validation, relying on toy examples or author-run case studies rather than
  independent evaluations with real teams in production
  environments~\citep{souza2019}.
\end{enumerate}

The mapping study also reported that 76.9\% of the 39 primary studies addressed functional
requirements, with a significant proportion also treating non-functional requirements;
yet the methods for handling quality attributes — precisely the domain where architectural
uncertainty is most consequential — were among the least well supported by tools or
empirical evidence \citep{souza2019}.

These findings reinforce the motivation for this work on two fronts. First, they confirm
the existence of a ``tooling void'' in which architectural management remains
disconnected from the daily engineering workflow. Second, they underscore the need for
empirical evaluations with real teams -- precisely the kind of evaluation this work aims
to provide for the Hypo-Stage tool and ArchHypo.
